\SourceLeaf{91}
$u$ ipso $t$ minorem fore. Denique $-2$ etiam ipsius $u$ residuum erit, \emph{i. e.} $t$ non erit minimus numerus qui inductioni nostrae adversatur, contra hyp. Quare necessario $-2$ omnium numerorum formarum $8n+5$, $8n+7$ non-residuum.

Combinando haec cum prop. art. 111, prodeunt theoremata haec:

I.\quad \emph{Omnium numerorum primorum $8n+5$, tum $-2$ tum $+2$ sunt non-residua,} uti iam in art. praec. invenimus.

II.\quad \emph{Omnium numerorum primorum formae $8n+7$, $-2$ est non-residuum, $+2$ vero residuum.}

Ceterum in utraque demonstratione pro $a$ etiam valorem parem accipere potuissemus; tunc autem casum ubi $a$ fuisset formae $4n+2$, ab eo distinguere oportuisset, ubi $a$ formae $4n$. Evolutio autem perinde procedit uti supra, nullique difficultati est obnoxia.

\Article{114}
Unus adhuc superest casus, scilicet ubi numerus primus est formae $8n+1$
Hic vero methodum praecedentem eludit, artificiaque prorsus peculiaria postulat.

Sit pro modulo primo $8n+1$, radix quaecunque primitiva $a$, eritque (art. 62) $a^{4n}\equiv-1\,(\operatorname{mod.}8n+1)$, quae congruentia ita etiam exhiberi potest, $(a^{2n}+1)^2\equiv2a^{2n}\,(\operatorname{mod.}8n+1)$, sive etiam ita, $(a^{2n}-1)^2\equiv-2a^{2n}$. Unde sequitur tum $2a^{2n}$ tum $-2a^{2n}$ ipsius $8n+1$ esse residuum: at quia $a^{2n}$ est quadratum per modulum non divisibile, manifesto etiam tum $+2$ tum $-2$ residua erunt (art. 98).

\Article{115}
Haud inutile erit, adhuc aliam huius theorematis demonstrationem adiicere, quae similem relationem ad praecedentem habet, ut theorematis art. 108 demonstratio secunda (art. 109) ad primam (art. 108). Periti facilius tunc perspicient, binas demonstrationes tam illas quam has non adeo heterogeneas esse, quam primo forsan aspectu videantur.

I.\quad Pro modulo quocunque primo formae $4m+1$, inter numeros ipso minores $1,2,3\ldots4m$, reperientur $m$ qui biquadrato congrui esse possunt. reliqui vero $3m$ non poterunt.

Facile quidem hoc ex principiis Sect. praec. derivatur, sed etiam absque his demonstratio haud difficilis. Demonstravimus enim pro tali modulo $-1$ sem\SourceLeaf{92}per esse residuum quadraticum. Sit itaque $ff\equiv-1$ patetque, si $z$ fuerit numerus quicunque per modulum non divisibilis, quaternorum numerorum $+z,-z,+fz,-fz$ (quos incongruos esse facile perspicitur) biquadrata inter se congrua fore; porro manifestum est biquadratum numeri cuiuscunque, qui nulli ex his quatuor congruus. illorum biquadratis congruum fieri non posse, (alias enim congruentia $x^4\equiv z^4$ quae est quarti gradus plures quam $4$ radices haberet, contra art. 43). Hinc facile colligitur, omnes numeros $1,2,3,\ldots4m$, tantummodo $m$ biquadrata incongrua praebere, quibus inter eosdem numeros $m$ congrui reperientur, reliqui autem nulli biquadrato congrui esse poterunt.

II.\quad Secundum modulum primum formae $8n+1$, $-1$ biquadrato congruus fieri poterit ($-1$ erit \emph{residuum biquadraticum} huius numeri primi).

Omnium enim residuorum biquadraticorum ipso $8n+1$ minorum (cifra exclusa) multitudo erit $=2n$ \emph{i. e.} par. Porro facile probatur, si $r$ fuerit residuum biquadraticum ipsius $8n+1$, etiam valorem expr. $\frac{1}{r}\,(\operatorname{mod.}8n+1)$ fore tale residuum. Hinc omnia residua biquadratica in classes simili modo distribui poterunt, uti in art. 109 residua quadratica distribuimus: nec non reliqua demonstrationis pars prorsus eodem modo procedit ut illic.

III.\quad Iam sit $g^4\equiv-1$, et $h$ valor expr. $\frac{1}{g}\,(\operatorname{mod.}8n+1)$. Tunc erit
\[
(g\pm h)^2=g^2+h^2\pm2gh\equiv g^2+h^2\pm2
\]
(propter $gh\equiv1$). At $g^4\equiv-1$, adeoque $-h^2\equiv g^4h^2\equiv g^2$, unde tandem $g^2+h^2\equiv0$, atque $(g\pm h)^2\equiv\pm2$ \emph{i. e.} tum $+2$ tum $-2$ residuum quadraticum ipsius $8n+1$. \emph{Q. E. D.}

\Article{116}
Ceterum ex praecc. facile regula sequens generalis deducitur: $+2$ \emph{est residuum numeri cuiusvis, qui neque per $4$, neque per ullum primum formae $8n+3$ vel $8n+5$ dividi potest, reliquorum autem} (\emph{ex. gr.} omnium numerorum formarum $8n+3$, $8n+5$, sive sint primi, sive compositi) \emph{non-residuum.}

$-2$ \emph{est residuum numeri cuiusvis, qui neque per $4$, neque per ullum primum formae $8n+5$ vel $8n+7$ dividi potest, omnium autem reliquorum non-residuum.}

Theoremata haec elegantia iam sagaci Fermatio innotuerunt, \emph{Op. Mathem. p.} 168. Demonstrationem vero quam se habere professus est, nusquam commu\SourceLeaf{93}nicavit. Postea ab ill. Eulero frustra semper est investigata: at ill. La Grange primus demonstrationem rigorosam reperit, \emph{Nouv. M\'em. de l'Ac. de Berlin} 1775. \emph{p.} 349, 351. Quod ill. Eulerum adhuc latuisse videtur, quando scripsit diss. in \emph{Opusc. Analyt.} conservatam, T. I. p. 259.

\begin{center}\emph{Residua $+3$ et $-3$.}\end{center}

\Article{117}
Pergimus ad residua $+3$ et $-3$. A posteriori initium faciamus.

Reperiuntur ex tab. II. numeri primi quorum residuum est $-3$ hi: 3, 7. 13, 19, 31, 37, 43, 61, 67, 73, 79, 97, inter quos nullus invenitur formae $6n+5$. Quod vero etiam ultra tabulae limites nulli primi huius formae dantur quorum residuum $-3$, ita demonstramus: Primo patet quemvis numerum compositum formae $6n+5$ necessario factorem primum aliquem eiusdem formae involvere. Quousque igitur nulli numeri primi formae $6n+5$ dantur, quorum residuum $-3$, eousque tales etiam compositi non dabuntur. Quodsi vero ultra tabulae nostrae limites tales numeri darentur, sit omnium minimus $=t$, ponaturque $-3=aa-tu$. Tunc erit, si acceperis $a$ parem ipsoque $t$ minorem, $u<t$, atque $-3$ residuum ipsius $u$. Sed quando $a$ formae $6n\pm2$, $tu$ erit formae $6n+1$, adeoque $u$ formae $6n+5$. \emph{Q. E. A.} quia $t$ minimum esse numerum inductioni nostrae adversantem supposuimus. Quando vero $a$ formae $6n$, erit $tu$ formae $36n+3$ adeoque $\frac{1}{3}tu$ formae $12n+1$, quare $\frac{1}{3}u$ erit formae $6n+5$; patet autem $-3$ etiam ipsius $\frac{1}{3}u$ residuum fore, atque esse $\frac{1}{3}u<t$, \emph{Q. E. A.} Manifestum itaque, $-3$ nullius numeri formae $6n+5$ residuum esse posse.

Quoniam quisque numerus formae $6n+5$ necessario vel sub forma $12n+5$. vel sub hac $12n+11$ continetur, prior autem forma sub hac $4n+1$, posterior sub hac $4n+3$, haec habentur theoremata:

I.\quad \emph{Cuiusvis numeri primi formae $12n+5$, tum $-3$ tum $+3$ non-residuum est.}

II.\quad \emph{Cuiusvis numeri primi formae $12n+11$, $-3$ est non-residuum, $+3$ vero residuum.}

\SourceLeaf{94}
\Article{118}
Numeri quorum residuum est $+3$ ex tabula II. inveniuntur hi: 3, 11, 13, 23, 37, 47, 59, 61, 71, 73, 83, 97, inter quos nulli sunt formae $12n+5$ vel $12n+7$. Nullos autem omnino numeros formarum $12n+5$, $12n+7$ dari quorum $+3$ sit residuum, eodem prorsus modo, ut in artt. 112, 113, 117, comprobari potest, quare hoc negotio supersedemus. Habemus itaque collato art. 111 theoremata:

I.\quad \emph{Numeri cuiusvis primi formae $12n+5$ non-residua sunt tum $+3$ tum $-3$} (uti iam in art praec. invenimus).

II.\quad \emph{Numeri cuiusvis primi formae $12n+7$ non-residuum est $+3$, $-3$ vero residuum.}

\Article{119}
Nihil autem per hanc methodum pro numeris formae $12n+1$ inveniri potest, qui proin artificia singularia requirunt. Ex inductione quidem facile colligitur, omnium numerorum primorum huius formae residua esse $+3$ et $-3$. Manifesto autem demonstrari tantummodo debet, numerorum talium residuum esse $-3$, quia tunc necessario etiam $+3$ residuum esse debet (art. 111). Ostendemus autem generalius, $-3$ esse residuum numeri cuiusvis primi formae $3n+1$.

Sit $p$ huiusmodi primus atque $a$ numerus pro modulo $p$ ad exponentem $3$ pertinens (quales dari ex art. 54 manifestum, quia $3$ submultiplum ipsius $p-1$). Erit itaque $a^3\equiv1\,(\operatorname{mod.}p)$ \emph{i. e.} $a^3-1$ sive $(a^2+a+1)(a-1)$ per $p$ divisibilis. Sed patet $a$ esse non posse $\equiv1\,(\operatorname{mod.}p)$, quia $1$ ad exponentem $1$ pertinet, quare $a-1$ per $p$ divisibilis non erit, sed $a^2+a+1$ erit, hincque etiam $4aa+4a+4$, \emph{i. e.} erit $(2a+1)^2\equiv-3\,(\operatorname{mod.}p)$ sive $-3$ residuum ipsius $p$. \emph{Q. E. D.}

Ceterum patet, hanc demonstrationem (quae a praecedentibus est independens) etiam numeros primos formae $12n+7$ complecti quos iam in art. praec. absolvimus.

Observare adhuc convenit, hanc analysin ad instar methodi in artt. 109, 115 usitatae exhiberi posse, at brevitatis gratia huic rei non immoramur.
